% 15_reference_impl.tex -- appendix
\section{Reference Detector Implementation}
\label{app:refimpl}

The following 24-line Python listing is a self-contained,
dependency-minimal reference implementation of
\autoref{alg:detect} for clockwork MCL. It takes a list of integer
token IDs, a bytes key, a state count $S$, and a significance
threshold $\alpha$; it returns the fingerprint score $\phi$, the
$z$-score, the one-sided $p$-value, and a binary detection flag.

The listing is sufficient to independently verify any MCL-watermarked
text given only the secret key, with no language-model access and no
learned components. Extensions to arbitrary $k$-regular topologies
(cf.~\autoref{prop:kregular}) require replacing the successor check
\texttt{states[i+1] == (states[i]+1) \% S} with a lookup
\texttt{T[states[i]][states[i+1]] > 0} and setting \texttt{p0 =
(T > 0).mean()}; no other modification is needed.

\begin{figure*}[t]
\begin{verbatim}
import hashlib
from math import sqrt
from scipy.stats import norm

def sha_state(key: bytes, token_id: int, S: int) -> int:
    """SHA-256-based state assignment sigma_kappa(t) = H(kappa || t) mod S."""
    h = hashlib.sha256(key + token_id.to_bytes(8, "big")).digest()
    return int.from_bytes(h[:8], "big") % S

def detect_mcl(token_ids, key: bytes, S: int,
               alpha: float = 0.01) -> dict:
    """Clockwork MCL detector; O(n) in the number of tokens."""
    n = len(token_ids)
    if n < 2:
        return {"phi": 0.0, "z": 0.0, "p_value": 1.0,
                "is_watermarked": False}
    states = [sha_state(key, t, S) for t in token_ids]
    valid = sum(1 for i in range(n - 1)
                if states[i + 1] == (states[i] + 1) % S)
    phi = valid / (n - 1)
    p0 = 1.0 / S
    se = sqrt(p0 * (1.0 - p0) / (n - 1))
    z = (phi - p0) / se
    p_value = 1.0 - norm.cdf(z)
    return {"phi": phi, "z": z, "p_value": p_value,
            "is_watermarked": p_value < alpha}
\end{verbatim}
\caption{Reference detector for clockwork MCL: 24 lines of
  dependency-minimal Python (\texttt{hashlib} + \texttt{scipy.stats}).
  Inputs are the tokenizer's token IDs, the secret key, and the state
  count; outputs are the fingerprint $\phi$, $z$-score, one-sided
  $p$-value, and a binary detection flag.}
\label{fig:refimpl}
\end{figure*}

\paragraph{Complexity.} Two SHA-256 evaluations per token
($2n \cdot 256$ bits of hash output) and $n-1$ integer equality
checks. At $n = 100$ the entire detection pipeline runs in under
$1$~ms on a single CPU core; no GPU, no model weights, no tokenizer
aside from what is needed to obtain \texttt{token\_ids}.

\paragraph{Interpretation as an audit primitive.} The detector's
inputs are exactly what a third-party auditor could receive under an
Article 50 disclosure regime: a public text, a detector program, and
a (confidentially held) key. The output is a $p$-value with
analytically known false-positive behaviour under the null,
sidestepping the calibration-by-grid-search problem that afflicts
logit-bias watermark families.

\paragraph{Key management.} The scheme's security reduces to the
confidentiality of \texttt{key} plus the random-oracle idealisation
of SHA-256 (\autoref{thm:security}). In practice an auditor and a
deployer can share the key through any standard key-management
substrate (e.g., HKDF-derived per-deployment keys, committed to an
external ledger so that the commitment precedes generation). Rotating
keys per deployment epoch preserves the $p_0 = k/S$ null baseline
without changing any detector behaviour.
